\section{Discussion}\label{sec:discussion}

\subsection{Answers to the research questions}

The Transformer autoencoder compresses variable-length PER careers well (RQ1): its held-out error is a third of that of a player-specific linear trend, and from $d=32$ the code is close to lossless. That success is exactly what makes the representation unsuitable for typing careers. A near-lossless code retains every idiosyncrasy of a career, and the geometry in which those idiosyncrasies are arranged is not identified by the reconstruction objective: two encoders trained from different seeds reconstruct equally well but place players differently relative to one another. Every downstream irreproducibility documented here, from the seed-dependence of single UMAP+HDBSCAN runs (RQ2) to the weak five-seed consensus (RQ3) and the low agreement across latent dimensionalities, follows from this property. The result is consistent with the comparative evidence that deep representations do not dominate classical ones in time-series clustering \citep{lafabregue2022} and adds a mechanism: without a clustering-aware objective \citep{ma2019} or a constraint on the latent geometry, reconstruction leaves the arrangement of the code under-determined, and structure discovery inherits that freedom as noise.

The comparison with simple representations (RQ4) is unambiguous. Nine summary statistics, or a ten-point resampling of the trajectory, yield partitions that are reproducible under subsampling, explain more trajectory-shape variance and are at least as strongly related to external outcomes. For univariate sequences of at most 23 points, hand-crafted features are not a weaker alternative to representation learning; they are the better tool.

The most consequential finding concerns the premise of the study. Reproducibility is necessary for a typology but not sufficient for the existence of types. Different reproducible representations partition the same players differently; the gap statistic finds no more structure than a single correlated cloud; the dip test finds no multimodality; and \nullDiscussion{} The picture is a continuum of career productivity paths, organised by level and length with secondary variation in direction, on which $k$-means draws reproducible but conventional boundaries. Three continuous coordinates, level, tilt and length, summarise that continuum with far less information loss than any typology and without the false precision of a boundary; we recommend them, or a functional-data equivalent \citep{wakim2014,elijah2025}, as the default description of a career in data of this kind, with types reserved for communication. The earlier claim of nineteen discrete types, and by implication similar claims obtained with single-run density-based pipelines, should be read in this light.

\subsection{Relation to prior work and hypotheses}

The coarse structure recovered here is compatible with the three peak-timing groups of \citet{wakim2014}, which the shape-only analysis reproduces almost literally (Table~\ref{tab:shape}); \citet{wakim2014} likewise found position unrelated to trajectory group. Our tests add that peak timing is continuously distributed and confounded with career length once trajectories are placed on normalised time. The findings are equally compatible with aging-curve studies that treat individual departures from the population curve as continuous random effects \citep{vaci2019,elijah2025}; the random-effects view is the more accurate description for this population, and typologies are best treated as descriptive discretisations, as in group-based trajectory modelling when class separation is weak \citep{nagin2005}. The association of the types with career length and honours echoes the retention literature: consistently below-average productivity shortens careers \citep{deutscher2017,coates2010}, and the short-career types are dominated by second-round and undrafted players \citep{berri2011draft,motomura2016}.

Several features of the profiles invite explanation that the data cannot provide, and we state them as hypotheses. The early-peak, medium-length type (T2) is where careers interrupted by serious injury, or altered by a change of role, would be expected to appear \citep{drakos2010,harris2013,defroda2021}; the late-rising type (T4) is consistent with players who obtained a larger role late; the persistence of a low-PER, long-career type (T6) with almost no honours is consistent with the value of skills that PER does not measure \citep{kubatko2007,berri2010} and with the incentive created by pension vesting after three seasons \citep{nbpa2023}. Whether injury, opportunity, role or unmeasured skill account for membership is a question for studies that include those variables.

\subsection{Methodological implications and use of the typology}

Three lessons follow for structure discovery in sports analytics. A single run of an embedding-plus-density pipeline should not be reported without a seed and subsample stability analysis; a mean ARI of at least 0.8 is a reasonable minimum. Internal indices computed in the clustering space cannot arbitrate between representations, and rules that mix several indices with ad hoc weights can pick out almost any solution: the three rules in Appendix Table~\ref{tab:selection} selected three different solutions from the same grid. And reproducibility must be checked against a structureless reference before a partition is called a typology. The eight types themselves are a stable, interpretable segmentation of complete careers by level, length and direction of PER, useful as a descriptive vocabulary and as strata in studies of career outcomes. They carry no evidence of predictive validity: they are defined from complete careers, no prospective evaluation was performed, and nothing here supports their use in scouting or roster decisions; that would require typing careers from their first seasons and testing the prediction of later outcomes on held-out cohorts.
